%%=============================================================================
%% CAPITULO 2 -- REFERENCIAL TEORICO
%% Texto de esqueleto (provisorio), a ser reescrito pela autora.
%%=============================================================================

\section{Referencial teórico}

O referencial teórico desta pesquisa parte da Lei Geral de Proteção de Dados
Pessoais, que estabelece regras para o tratamento de dados pessoais, inclusive
pelo poder público \parencite{brasil2018}, e a relaciona com o uso de
inteligência artificial na gestão documental realizada no Sistema Eletrônico de
Informações. Sobre as decisões automatizadas, Almada e Maranhão registram:

\begin{citacao}
Sistemas de tomada de decisão automatizada são cada vez mais comuns em diversas
aplicações nos setores público e privado. Como muitas das decisões tomadas por
estes sistemas dependem de dados pessoais como insumos ou geram dados
associáveis a pessoa natural, a Lei Geral de Proteção de Dados (LGPD) introduz
regras específicas voltadas a decisões baseadas unicamente no tratamento
automatizado de dados pessoais. \parencite{almada2023}
\end{citacao}

No campo da conformidade, \citeonline{dias2023} discutem os riscos e os
benefícios do uso da inteligência artificial na atividade de compliance,
discussão que orienta a proposta de monitoramento desta pesquisa. A
Figura~\ref{fig:fluxo-monitoramento} apresenta o fluxo pretendido, dos
documentos do SEI ao alerta de conformidade, e a
Tabela~\ref{tab:categorias-dados} reúne as categorias de dados a monitorar.

\begin{figure}[htb]
  \caption{Fluxo do monitoramento de conformidade no SEI}
  \label{fig:fluxo-monitoramento}
  \centering
  \begin{tikzpicture}[
      node distance=0.6cm,
      etapa/.style={draw, rounded corners, align=center, font=\small,
                    text width=2.6cm, minimum height=1.4cm},
      seta/.style={-{Stealth}, thick}]
    \node[etapa] (sei) {Documentos do SEI};
    \node[etapa, right=of sei] (pln) {Processamento de linguagem natural};
    \node[etapa, right=of pln] (dados) {Identificação de dados pessoais};
    \node[etapa, right=of dados] (alerta) {Alerta de conformidade};
    \draw[seta] (sei) -- (pln);
    \draw[seta] (pln) -- (dados);
    \draw[seta] (dados) -- (alerta);
  \end{tikzpicture}
  \fonte{Elaborado pela autora (2026).}
\end{figure}

\begin{table}[htb]
  \caption{Categorias de dados a monitorar}
  \label{tab:categorias-dados}
  \centering
  \begin{tabular}{ll}
    \toprule
    \textbf{Categoria} & \textbf{Exemplos} \\
    \midrule
    Dados pessoais           & nome, CPF, endereço, telefone \\
    Dados pessoais sensíveis & saúde, religião, biometria \\
    Dados de crianças e adolescentes & nome e data de nascimento de menores \\
    \bottomrule
  \end{tabular}
  \fonte{Elaborado pela autora (2026).}
\end{table}
